% ============================================================
%  H. Høffding, "Om Vitalisme"
%  A lecture delivered in the Biological Society, 27 January 1898.
%  Printed in: Mindre Arbejder (1899), pp. 40–50.
%  TRANSLATION (English)
%
%  Base text: the Danish of Mindre Arbejder as transcribed on Danish
%  Wikisource (da.wikisource.org/wiki/Om_Vitalisme), transcluded from the
%  page scan Indeks:Høffding - Mindre Arbejder.djvu, scan pages 54–64 =
%  printed pp. 40–50. 253 of that scan's 262 pages carry the quality level
%  *Valideret* — proofread by two independent users against the images —
%  and every page of this essay is among them. Public domain: Høffding died
%  1931; first published 1899.
%
%  Printed pages are marked in this file as % [p. N] comments at the point
%  where the page turns, so the translation can be collated against the
%  original without cluttering the output.
%
%  READINGS QUERIED IN THE BASE TEXT — see transcription.tex for the full
%  log. One is substantive: at printed p. 48 the Wikisource text reads
%  "den Forandring, hvis teoretiske Udtryk den gamle Vitalisme var"
%  (the *change* whose theoretical expression the old vitalism was), which
%  makes no sense in context and is almost certainly "Forundring" —
%  *wonder*, *astonishment* — misread a-for-u. The following sentence
%  confirms it: "Nyvitalismen er ogsaa Udtryk for en For[u]ndring … der
%  vækkes ved det Hemmelighedsfulde." Translated here as "wonder".
% ============================================================
\documentclass[12pt,a4paper]{article}
\usepackage[utf8]{inputenc}
\usepackage[T1]{fontenc}
\usepackage{libertinus}
\usepackage{libertinust1math}
\usepackage[english]{babel}
\usepackage[protrusion=true,expansion=true]{microtype}
\usepackage{setspace}
\usepackage[margin=1.2in]{geometry}
\usepackage{fancyhdr}
\usepackage{hyperref}

\hypersetup{
  pdftitle={On Vitalism},
  pdfauthor={H. Høffding},
  hidelinks
}

\pagestyle{fancy}
\fancyhf{}
\rhead{Høffding, \emph{On Vitalism} (1898)}
\cfoot{\thepage}

\setstretch{1.3}

\title{\textbf{On Vitalism}\footnote{A lecture delivered in the Biological
Society, 27 January 1898.}}
\author{H.\ Høffding\\[4pt]
  \small Translated from the Danish}
\date{1898}

\begin{document}
\maketitle
\thispagestyle{fancy}

\bigskip

% [p. 40]
\noindent It is more than a figurative expression when one speaks of thoughts
living and dying. Like all organic beings, thoughts have their determinate
conditions, on which their arising, their development or their disappearance
depends. Thoughts must, like organisms, struggle in order to live. Only those
thoughts live in the long run which can subsist with, and be the expression of,
large circles of well-tested and coherent experience. And just as a natural
selection takes place among the organic germs, tentative beginnings and
variations which are present at every point of development, so too a choice is
constantly made among the thoughts and points of view that are available at any
given time. There is a reciprocal relation between thinking and experience. For
the involuntarily formed thoughts and points of view among which experience
chooses have of course not themselves arisen without the influence of
experience. But just as in nature far more germs and variations arise than are
selected for continued development, so thinking too begins with energetic and
luxuriant sallies and anticipations, which are gradually sifted, limited,
developed or discarded.

This holds especially as regards the understanding of organic life. The struggle
between speculation and verification moves in rhythmic fashion through the
history of biology, and in the most recent period a swell seems to be astir
which is leading to the taking up again of points of view whose time seemed
decidedly to be past. Life has been thought to be found in old thoughts
% [p. 41]
whose decease had already been announced. In what sense that opinion is
justified, we shall examine more closely.

There are two fundamental conceptions which here stand over against one another.
The one rests on this, that the clearest understanding we can have of a
phenomenon is the one we have where we can resolve it into its elements and see
it again as the product or result of those elements' co-operation --- that is,
when we go from the parts to the whole. The so-called mechanical view of nature
maintains that this mode of explanation holds not only for the physical and
chemical phenomena but for all natural phenomena. Against this it is urged from
the other side that what is peculiar to the organic phenomena is that the
constitution, the form and the mode of operation of the elements, of the parts,
cannot be understood except by taking as one's ground the thought of the whole
which they form, or are to form. The organism cannot be explained as a mere
product of organs or cells, nor the cell as a mere product of molecules. We must
then here assume a peculiar force, operating according to other laws than the
mechanical ones. This is the vitalistic conception. --- And since the only way in
which we can represent a whole as determining the constitution and the
co-operation of the parts is to represent it as a purpose that is striven after,
a vitalistic conception will always remain more or less teleological.

When history shows us a rhythmic alternation of these opposed ways of seeing,
the advance of knowledge is by no means thereby excluded. It is the new
experiences that let the phenomena appear now more from the one side, now more
from the other, and that also gradually transform the earlier forms of the two
fundamental conceptions.

A mechanical view of nature first comes forward at a relatively advanced
standpoint. Primitive man's conception is vitalistic or animistic, explaining
all phenomena that awaken attention by the intervention of personal beings. All
later vitalism has a certain kinship with this primitive conception. --- In
Greek antiquity the standpoint of mechanism is urged by the \emph{Atomists}
(whose doctrine was later set forth in Lucretius' famous poem \textit{de rerum
natura}).
% [p. 42]
\emph{Aristotle}, who was himself a brilliant investigator of nature in the
domain of organic life, maintained, in opposition to atomism, a conception which
regarded nature as a great scale of stages, where every member is matter and
means in relation to a higher one, while it is form and purpose in relation to a
lower. Nowhere in nature could he conceive the whole as arising as the result of
the elements, the form as the effect of the matter. He therefore extends
vitalism to the whole of nature. It is thus characteristic of the position of
the problem of life in antiquity that no sharp boundary was set between
inorganic and organic phenomena. For the Aristotelian conception as for the
atomistic there are in nature only differences of degree.

For the time being the Aristotelian conception was victorious, and it became the
prevailing one in the Middle Ages. A turning point comes with \emph{the founding
of modern natural science} by Leonardo da Vinci, Kepler and Galilei. The
proposition was now set up that all change is motion, and it was demanded that
this be carried through in every domain, so that all laws of nature might be
derived from the principles of mechanics. \emph{Cartesian philosophy}
generalized this conception and carried it out into wide circles. The discovery
of the circulation of the blood and of reflex movement was an empirical support
for the conviction that the domain of organic life could form no exception.

As a reaction against this whole tendency there arises the conception one now
chiefly thinks of when one speaks of vitalism. \emph{G. E. Stahl} combated
mechanical physiology and maintained (especially in \textit{Theoria medica vera}
1708) that without a soul's intervention organic growth and the organism's
self-preservation were inexplicable. He thus goes further back than to
Aristotle, in that he renews the primitive animism. He held that since the
conscious soul can be the cause of voluntary movements, it must also be able to
be the cause of the involuntary processes in the organism, even if it here
operates in a more instinctive manner, and even if it does not
% [p. 43]
always behave rationally --- which, after all, it does not always do in its
voluntary actions either. --- The well-known medical school at Montpellier
modified the theory by distinguishing between soul and life, so that between
soul and body a peculiar life-principle was inserted. In this form vitalism is
presented by \emph{Barthez} (\textit{Nouveaux Elements de la science de
l'homme.} Montpellier 1778). He combats two sects, the Mechanists and the
Animists, and represents vitalism in the narrowest sense of the word. He defines
the life-principle as ``the cause which brings forth all the phenomena of life
in the human body''. It follows laws that lie above the laws of physics
(\textit{d'un ordre transscendant par rapport aux lois de la Physique}). Its
proper nature, and especially its relation to the soul proper, is unknown. But
Barthez plainly enough conceives it on the analogy of mental life proper; he
ascribes to it instinct, temperament and sensory capacity. It intervenes in the
life-process and directs it in a determinate direction. It lays in the muscles
the force needed to overcome a determinate resistance, and shows its power
especially in convulsive attacks. It dilates the muscle fibres in the ducts of
milk and of semen, and it acts upon the humours, so that the substances are
formed which are needed for repair or for reproduction. It regulates respiration
and the development of heat, and conditions a ``sympathetic'' connection between
the organism's different parts. --- Different both from Barthez's vitalism and
from Stahl's is the theory which the famous anatomist \emph{Bichat} set up (in
\textit{Recherches physiologiques sur la vie et la mort.} Paris 1800). He
conceives the vital force as a property of the organic tissues. It was an
incipient decentralization of vitalism.

\emph{Classical vitalism} --- so I designate the conception represented, in
different shadings, by Stahl, Barthez and Bichat --- has this in common with the
Aristotelian, that by assuming a vital force it takes itself to have given a
sufficient explanation of the peculiarity of the phenomena of life. It differs
from Aristotelian vitalism in sharply emphasizing the opposition between the
inorganic and the organic, and even assuming a hostile relation
% [p. 44]
between them. The soul, or the vital force, must struggle with the mechanical
and chemical conditions in order to get them to operate in a determinate
direction; indeed Bichat even defines life as the sum of the forces that
struggle against death.

A new and magnificent rise of the mechanical view of nature in its application
to biology came about on the basis of the work of Italian, French and Swiss
investigators of nature at the close of the eighteenth century --- work which,
however, only came to be properly recognized towards the middle of our own
century. About the year 1840 \emph{Dumas} and \emph{Liebig} reformed plant
physiology and demonstrated the circulation of matter in nature, and with it the
connection between organic and inorganic nature. Soon after, \emph{Robert Mayer}
discovered the law of the conservation of energy and applied it at once to the
organic domain (in his work \textit{Die organische Bewegung} 1845). A
proposition from that work may stand as the motto of the whole new biological
conception: ``In the process of life there takes place only a transformation of
matter and of force, but never a creation of the one or the other.'' A whole
series of eminent physiologists, German ones above all (notably the disciples of
Johannes Müller --- though the master himself was still a vitalist), showed the
fruitfulness of the new points of view. Physiology now seemed able to be
conceived as applied physics and chemistry. The assumption of a vital force was
stamped as mysticism and superstition. \emph{Helmholtz} uses (in his work
\textit{Das Denken in der Medicin} 1878) vitalism as the opposite of natural
inquiry, and regards the question of the vital force as one with the question of
the perpetuum mobile. And in his memorial address on Helmholtz (1897)
\emph{Dubois Reymond} declared that the doctrine of the conservation of energy
gave the phantom of a vital force its final blow. The period in biology which
began about 1840 is one of the most fruitful in its whole history, especially
through its many applications of physical and chemical laws to the organic
domain. A great and coherent view of nature might hereby seem to have been won
once for all.

In the latest years one has begun to speak again of vitalism as a necessary
standpoint. Connected with this in a certain way is the fact that one speaks of
natural science having
% [p. 45]
disappointed the expectations that had been placed in it --- indeed, that one
has even spoken of ``the bankruptcy of science''. The facts on which so-called
\emph{neo-vitalism} rests were in the main already known earlier, and were
emphasized by critical investigators. Men like Lotze and Virchow, Littré and
Claude Bernard, who all maintained that only the methods of physics and
chemistry could lead to insight into the laws of vital phenomena, nevertheless
by no means held that life's innermost nature could be fully understood from
those laws alone. They regarded life as a fact, and saw it as biology's task to
investigate the single members of the life-process as it actually runs. The
occasion of neo-vitalism's appearance can be sought only in this, that many
parts of the life-process have shown themselves to be \emph{still} more
intricate than had been thought beforehand. There are three points especially
that come into the foreground for neo-vitalism.

1) Even if at every single transition in the life-process there occurred only
such transformations of matter and force as could be explained by chemical and
physical laws, there would still remain the fact that the chemical and physical
processes in growth and nutrition run in determinate \emph{directions}. This is
all the more striking in that growth takes place sporadically, from different
points, so that the joining into a form of the whole comes about only gradually.
Should it not then be justified to speak of a formative drive (\textit{nisus
formativus})?

2) In \emph{the taking up of nutriment} into the cells a choice is exercised.
Neither the quantity nor the quality of what is taken up can be explained in a
purely chemical or physical way. The cells seem here to act as little beings
whose need and wish determine their relation to the surrounding world. Several
neo-vitalists therefore personify the cells, as the classical vitalists
personified the whole organism. Since we know only from ourselves what need and
wish are, one of the most zealous neo-vitalists, \emph{Bunge}, has even declared
that psychology ought properly, here, to come to physiology's aid.

3) Organic tissue differs from the crystal in that its molecules \emph{never}
come into \emph{complete equilibrium}
% [p. 46]
so long as life lasts. To be able to die is therefore a mark of the living. ---

It is beyond doubt that attention to these points has been sharpened in recent
years. It is doubtful, on the other hand, whether a new stage has really thereby
emerged in the old rhythmic process, such that the name neo-vitalism could be
justified. How far one will extend the name may in any case be doubtful.
Investigators like \emph{Giulio Fano} (\textit{La fisiologia come scienza
autonoma.} Genova 1884) and \emph{Hans Driesch} (\textit{Die Biologie als
selbständige Grundwissenschaft.} Leipzig 1893), who wish only to maintain
biology's independence over against chemistry and physics, can be called
neo-vitalists only improperly. The name does fit \emph{G. Bunge}, on the other
hand (the chapter ``Vitalismus und Mechanismus'' in ``Lehrbuch der
physiologischen und pathologischen Chemie''. Leipzig 1887), who almost renews
Stahl's animism. From standpoints very different from one another,
\emph{Haeckel}, \emph{Verworn} and \emph{Rindfleisch} have expressed themselves
in a similar direction.\footnote{The significant difference between the
conceptions that are often designated together as neo-vitalism was already
pointed out by \emph{Felice Tocco} in his essay ``Le disfatte della scienza''
(Nuova Anthologia. 1 March 1896).}

\bigskip

It will have its interest to compare neo-vitalism with classical vitalism, in
order to see whether it is, as some hold (for example \emph{Mosso} in his essay
``Materialismo e Misticismo'' in Nuova Anthologia, 1 December 1895), a sign of
retrogression, or whether, in relation to earlier and more or less kindred
tendencies, differences can be shown which point to advance. Such a comparison
can be led back to three principal points.

1) Classical vitalism holds that with the assumption of a ``vital force'' a
conclusive, positive explanation has been reached. But for the most critical
investigators, whom neo-vitalism claims for itself,
% [p. 47]
it is merely a matter of maintaining that biology has its independent empirical
foundation, and is something other and more than a sum of physical and chemical
deductions. In so far as a vital force is spoken of, the word means an $x$, an
unresolved remainder, and stands as the expression of the intricate conditions
under which the general laws of nature find application in the organic domain.

2) Classical vitalism was called by Claude Bernard ``a doctrine that makes
lazy'' (\textit{une doctrine paresseuse}). In its confidence in the vital force
it disdained exact investigation of the particular relations. Barthez held, for
example, that it was unnecessary and unreasonable to seek an anatomical basis
for the connection between sensory and motor functions; the unity of the vital
force was sufficient explanation of that connection. He thus rejected in advance
investigations such as those which later led Charles Bell to his discovery of
the difference between sensory and motor nerves. (See Barthez: \textit{Nouveaux
Elements}, p.~142--153; 207~f.)\footnote{This point was already brought out by
\emph{Louis Peisse} in his critique of the Montpellier school. (\textit{La
Médicine et les Medecins.} Paris 1857. I, p.~291--293).} Neo-vitalism, by
contrast, holds fast --- despite its inclination to psychological analogies ---
that the only method which can actually be in question for use in biology is the
one the mechanical conception has proclaimed. It regards the problem of life as
more intricate than classical vitalism did.

3) The most striking difference between neo-vitalism and classical vitalism
consists in this, that the vital force is not, or not only, ascribed to the
organism as a whole, but to the cells, the organism's elements. In place of a
genius directing the organism's life as a whole, one is now inclined to assume a
kind of cell-geniuses, which decide what and how much is to be taken up from the
inner medium. \emph{Bunge} says outright: ``\textit{Wir übertragen das aus dem
eigenen Bewusstsein geschöpfte auf die Objecte unserer Sinneswahrnehmung, auf
die Organe, auf die Gewebselemente, auf jede kleine Zelle.}'' It undoubtedly
marks an advance that the concept of force is used of the single co-operating
elements instead of
% [p. 48]
en bloc of the whole organism. The problem has become more localized; the
riddles have been pushed further back. To be sure, the great problem then
remains, how the life of the many cells is combined into the organism's
collective life --- and it was this above all that awakened the wonder whose
theoretical expression the old vitalism was. Neo-vitalism too is the expression
of a wonder, but especially of the one awakened by what is mysterious in the
processes that go on in the organism's smallest parts. One had not expected that
the problem of life would come again at the very elements into which one thought
one had resolved the whole of life. ---

The difference between neo-vitalism and classical vitalism is so great that it
is really unwarranted to say that it is one and the same conception that has
come again. With respect to what is scientifically asserted, and to the method
actually employed, no change has really taken place in relation to the
immediately preceding period in the history of biology. All talk of retrogression
(Mosso) or of bankruptcy (Brunetière) is therefore unwarranted, and the name
neo-vitalism itself had best be avoided.

\bigskip

In conclusion it will be of interest to bring out the analogy between the
problem which makes biology something other and more than applied chemistry and
physics, and a whole series of other problems which all rest on the opposition
between the rational and the empirical --- or, as one may also put it, between
the formal and the real in our knowledge. The more rational and formal our
knowledge is, the further it can move forward through definitions and
deductions, so that every transition takes place with intuitable necessity.
Empirical or real knowledge, by contrast, must often stop at facts which can
indeed be described and analysed, but not defined and derived from other facts.
Such a fact is life; and therefore biology belongs, and will perhaps always
belong, to empirical knowledge. It was not only Stahl who in his day missed a
definition of what life is (a lack which he did, however, believe himself able
to supply); it has been said in our own day as well (by Lotze and by Littré, by
Bernard and by Fano) that no such definition has been found. But life does not
stand alone in this respect. At a whole series of points
% [p. 49]
the rational continuity in our knowledge is replaced by empirical data which for
the time being stand as unresolvable. They are kinks in the thread, or perhaps
even knots.

1) Let us suppose that the proposition that all change is motion could be
completely carried through, and that the ideal of the mechanical view of nature
could so far be realized. The empirical kinks or knots in the thread would not
on that account be lacking. At every point --- however far back we go --- the
motion would still appear with a certain given direction and a certain given
velocity, and the outer causes which alter that direction and that velocity
would operate from certain determinate places and in certain determinate degrees
and directions. --- The sum of energy expended in the motion, and in the causes
that alter it, is moreover of a determinate given magnitude, and its distribution
among the single operative elements can likewise be known only empirically.

2) Energy appears, further, in qualitatively different forms, and the operative
elements likewise have their qualitative peculiarities. Even if, by the law of
the conservation of energy and of matter, equivalents can always be shown at
every transformation from one form of energy into another and at every change of
matter, there arise at the same time in these transformations new properties
which cannot be derived from the antecedent conditions, but are attached to them
purely empirically.

3) In the organic domain the sum of empirical elements becomes greater still,
and more pronounced. It is attention to these elements that often brings with it
a renewed inclination towards vitalistic views. As such elements we have in the
foregoing mentioned direction of growth and structure, qualitative and
quantitative choice in the cells' change of matter, and the maintenance of the
unstable equilibrium.

4) If we take the psychological domain along with the rest, the phenomena of
consciousness ---
% [p. 50]
from the most elementary sensation to the highest thoughts, feelings and acts of
will --- stand as an empirical given which cannot be derived from the material
presuppositions. Psychology is to a still higher degree than the other parts of
biology a purely descriptive science. And within psychology there then come
forward again new qualitative differences and transformations, where the same
problem presents itself anew.

It is clear that if, at all the points where the empirical, the qualitative and
the individual put to the test the rational continuity our knowledge strives
after, one were to assume new independent and original ``forces'', one would be
deceiving oneself. The limits we thus constantly come up against are always at
the same time to be regarded as new material, setting new tasks. It is
existence's richness and depth that produces these kinks and knots in the
thread, and thereby makes necessary a constant interaction between experience
and thinking, induction and deduction, analysis and synthesis --- that doubleness
in method which Goethe compared with the relation between breathing in and
breathing out. The struggle between mechanism and ``vitalism'' is, rightly
understood, a struggle between two tendencies between which there can at bottom
be no opposition. We should not be more perfect beings if we could breathe out
without breathing in. That doubleness makes both life and thought the richer.

\end{document}
